% Qualitative comparison (Section 3). Defensible, high-level claims only.
% Requires: booktabs. Keep textual (no pifont) for portability.
\begin{table*}[t]
  \centering
  \caption{Qualitative comparison of the evaluated frameworks. ``Own''
    means the capability is implemented in-framework without a
    third-party runtime library; ``Lib'' means it is delegated to a
    bundled external library. Security-by-default denotes whether the
    request pipeline applies protections (headers, CSRF, rate limiting,
    body caps) without opt-in configuration.}
  \label{tab:qualitative}
  \small
  \begin{tabular}{@{}lccccc@{}}
    \toprule
    & \textbf{JxMVC} & \textbf{Spring Boot} & \textbf{Quarkus} & \textbf{Micronaut} & \textbf{Javalin} \\
    \midrule
    Runtime dependencies   & \textbf{0}        & many        & many        & moderate    & few \\
    JSON                   & Own               & Lib (Jackson)& Lib        & Lib (Jackson)& Lib (Jackson) \\
    Connection pool        & Own (CAS)         & Lib (Hikari)& Lib (Agroal)& Lib (Hikari)& external \\
    Dependency injection   & Own (reflection)  & runtime     & compile-time& compile-time& none \\
    Security-by-default    & \textbf{Yes (pipeline)} & opt-in & opt-in     & opt-in      & manual \\
    OIDC (JDK-only)        & \textbf{Yes}      & Lib         & Lib         & Lib         & manual \\
    GraalVM native         & not yet           & partial (AOT)& first-class& first-class & limited \\
    Core artifact size     & \textbf{253\,KB}  & MB-scale    & MB-scale    & MB-scale    & MB-scale \\
    \bottomrule
  \end{tabular}
\end{table*}
